% Source: https://artofproblemsolving.com/community/c3h1717359_kootrapalis_easy_problems
% Stable identity: aops-topic-1717359-q05
Since $3^2+4^2=5^2$, the angle at $B$ is right. Put $B=(0,0)$, $C=(4,0)$, $A=(0,3)$, and let $\alpha=\angle BAC$, $\gamma=\angle BCA$, so $\tan\gamma=\dfrac{3}{4}$ and $\tan\alpha=\dfrac{4}{3}$. By the angle bisector theorem $\dfrac{BE}{EC}=\dfrac{AB}{AC}=\dfrac{3}{5}$, so $BE=\dfrac{3}{2}$ and $E=\left(\dfrac{3}{2},0\right)$.

\textbf{Locating $D$.} In the cyclic quadrilateral $CDFE$ the exterior angle at $D$ equals the interior angle at $E$: $\angle ADB=\angle ADF=\angle FEC=\angle AEC$. From triangles $ABD$ and $AEC$,
\[
180^\circ-\alpha-\angle ABD=180^\circ-\dfrac{\alpha}{2}-\gamma\quad\Longrightarrow\quad\angle ABD=\gamma-\dfrac{\alpha}{2} .
\]
Here $\tan\dfrac{\alpha}{2}=\dfrac{\sin\alpha}{1+\cos\alpha}=\dfrac{\ \dfrac{4}{5}\ }{\ \dfrac{8}{5}\ }=\dfrac{1}{2}$, so
\[
\tan\angle ABD=\dfrac{\dfrac{3}{4}-\dfrac{1}{2}}{1+\dfrac{3}{8}}=\dfrac{2}{11}.
\]
Since $\angle ABD$ is measured from $BA$ (the $y$-axis), line $BD$ is $x=\dfrac{2}{11}y$. It meets $AC$: $\dfrac {x}{4}+\dfrac {y}{3}=1$ at $D=\left(\dfrac{12}{25},\dfrac{66}{25}\right)$, and $AE$: $2x+y=3$ at $F=\left(\dfrac{2}{5},\dfrac{11}{5}\right)$.

\textbf{The centre.} $O$ lies on the perpendicular bisector $x=\dfrac{11}{4}$ of $CE$; write $O=\left(\dfrac{11}{4},k\right)$. From $OC^2=OD^2$:
\[
\dfrac{25}{16}+k^2=\left(\dfrac{227}{100}\right)^2+\left(k-\dfrac{66}{25}\right)^2\ \Longrightarrow\ \dfrac{132}{25}k=\dfrac{105600}{10000}\ \Longrightarrow\ k=2 .
\]
So $O=\left(\dfrac{11}{4},2\right)$ and the radius is $\sqrt{\dfrac{25}{16}+4}=\dfrac{\sqrt{89}}{4}$. (Check: $OF^2=\left(\dfrac{47}{20}\right)^2+\left(\dfrac{1}{5}\right)^2=\dfrac{2225}{400}=\dfrac{89}{16}$, so $F$ is on the circle too.)

\begin{asy}
usepackage("amsmath");
size(6cm);
defaultpen(fontsize(10pt));
pair A=(0,3),B=(0,0),C=(4,0),Ep=(1.5,0),D=(12/25,66/25),F=(0.4,2.2),O=(2.75,2);
draw(circle(O,sqrt(89)/4),blue+linewidth(0.6));
draw(A--B--C--cycle,linewidth(0.9));
draw(A--Ep,linewidth(0.7));draw(B--D,linewidth(0.7));
draw(O--A,dashed);draw(O--B,dashed);
draw(B--(0.18,0)--(0.18,0.18)--(0,0.18),linewidth(0.5));
dot(O);dot(F);dot(D);dot(Ep);
label("$A$",A,NW);label("$B$",B,SW);label("$C$",C,SE);label("$E$",Ep,S);
label("$D$",D,dir(-62)*2.6);label("$F$",F,W);label("$O$",O,E);
label("\mbox{}",(min(currentpicture,user=true)+max(currentpicture,user=true))/2,black); // keeps the text after the figure black when LaTeX typesets the labels inline
\end{asy}

Finally
\[
OA=\sqrt{\dfrac{121}{16}+1}=\dfrac{\sqrt{137}}{4},\qquad OB=\sqrt{\dfrac{121}{16}+4}=\dfrac{\sqrt{185}}{4},
\]
so $OA+OB=\dfrac{\sqrt{137}+\sqrt{185}}{4}$, and $a+b+c=137+185+4=\boxed{326}$.
